\chapter{Impact of Individual User Characteristics on Predictability}
\label{ch:user}

\section{Objectives and Predictability Determinants}

Identical fine-tuned models exhibit extreme performance variance across individual subscribers, yielding near 100\% prediction accuracy for certain individuals while scoring near zero for others. This variance indicates that aggregate population accuracy is incomplete without individual trajectory breakdown. The objective of this chapter is to resolve the following question: \emph{which intrinsic trajectory properties determine individual predictability?} If predictability depends primarily on hyperparameter tuning, model capacity should be increased. If predictability depends on individual subscriber mobility traits, aggregate performance is fundamentally constrained prior to training. The empirical results presented in this chapter demonstrate that individual subscriber characteristics dictate overall predictability.

\paragraph{Data Provenance.} Analysis uses two dataset cohorts from \textbf{Table~\ref{tab:datasets}}: raw daily logs evaluated on 1,173 subscribers, and active mobility logs (\Dmob{}) evaluating 400 training subscribers and \textbf{100 held-out subscribers}.

\section{Correlation Analysis of Individual User Features}

To identify trajectory features governing predictability, five subscriber-level candidate metrics computable directly from individual traces were evaluated: cell transition frequency, primary cell concentration, distinct cell count, transition temporal entropy, and total log event volume.

Each candidate metric was correlated with individual model prediction accuracy using the \textbf{Pearson correlation coefficient} ($r$), formulated according to Pearson \textbf{\cite{pearson1895correlation}}:

\begin{equation}
  r \;=\;
  \frac{\displaystyle\sum_{i=1}^{n} (x_i - \bar{x})\,(y_i - \bar{y})}
       {\sqrt{\displaystyle\sum_{i=1}^{n} (x_i - \bar{x})^2}\;
        \sqrt{\displaystyle\sum_{i=1}^{n} (y_i - \bar{y})^2}}
  \label{eq:pearson}
\end{equation}

where $n$ represents subscriber count, $x_i$ denotes candidate feature value for subscriber $i$, $y_i$ represents top-1 prediction accuracy for subscriber $i$, and $\bar{x}, \bar{y}$ denote sample means.

% F8 -- Where the evaluation signal lives, and what actually correlates with accuracy.
\begin{figure}[t]
\centering
\replegend{\swatch{clbase}~persistence baseline \qquad \swatch{clmodel}~model,
           \srawbig \qquad \swatch{clink}~model, \smob}
\fitwidth{%
\begin{tikzpicture}
\begin{groupplot}[
  group style={group size=1 by 2, vertical sep=1.9cm},
  rep, width=0.80\linewidth,
  title style={font=\footnotesize,yshift=1pt},
]
% ---- (a) volume quintiles ----
\nextgroupplot[repbar, height=4.6cm,
  title={(a) \srawbig --- five equal groups by predictions per user},
  symbolic x coords={va,vb,vc,vd,ve}, xtick=data,
  xticklabels={1--2,2--6,6--17,17--58,58--1{,}533},
  x tick label style={font=\scriptsize},
  xlabel={predictions per user}, ylabel={top-1 accuracy (\%)},
  ymin=0, ymax=100, ytick={0,20,...,100}, /pgf/bar width=9pt]
\addplot[fill=clbase,draw=none] coordinates {(va,66.3)(vb,61.9)(vc,62.2)(vd,66.6)(ve,78.1)};
\addplot[fill=clmodel,draw=none] coordinates {(va,66.0)(vb,61.9)(vc,61.9)(vd,68.1)(ve,78.6)};

% ---- (b) correlations ----
\nextgroupplot[rephbar, height=5.4cm,
  title={(b) Pearson correlation between a per-user feature and that user's accuracy},
  xmin=-1.2, xmax=1.2, xtick={-1,-0.5,0,0.5,1},
  xlabel={Pearson $r$}, /pgf/bar width=5.5pt,
  symbolic y coords={len,vol,uniq,irr,dom,ent,switch},
  ytick={len,vol,uniq,irr,dom,ent,switch},
  yticklabels={raw sequence length, number of predictions, unique cells visited,
               switch irregularity, dominant-cell freq., normalised entropy,
               \textbf{cell-switch rate}},
  y tick label style={font=\scriptsize,align=right},
  nodes near coords,
  /pgf/number format/.cd, fixed, fixed zerofill, precision=3, /pgfplots/.cd,
  every node near coord/.append style={font=\tiny}]
\addplot[fill=clmodel,draw=none] coordinates {(0.054,len)(0.154,vol)(-0.175,uniq)(0.696,dom)(-0.811,ent)(-0.971,switch)};
\addplot[fill=clink,draw=none]   coordinates {(-0.101,vol)(-0.232,uniq)(-0.140,irr)(0.285,dom)(-0.773,switch)};
\end{groupplot}
\end{tikzpicture}}
\caption[Concentration of the evaluation signal and per-user correlations]{(a)~The
evaluation signal is extremely concentrated: the median user contributes 11
predictions, while the most active tenth --- 117 people --- carries 58.6\,\% of
the 54{,}173 predictions. Accuracy nevertheless barely moves across the groups,
and the raw correlation with volume is weak ($r = +0.15$), so this is not a
volume effect: heavy users simply happen to switch cell less often
($r = -0.15$). (b)~Ranked by strength, the only feature that really explains
per-user accuracy is how often the user changes cell --- $r = -0.971$ on the
1{,}173-user set, confirmed at $r = -0.773$ on the 100 held-out users of a
different dataset. Features with a single bar were computed on one dataset only. Group sizes in
panel~(a), left to right: 324, 234, 234, 234 and 117 users, carrying 353, 948,
2{,}601, 7{,}553 and 42{,}718 predictions.}
\label{fig:concentration}
\end{figure}


\textbf{Figure~\ref{fig:concentration}} ranks candidate metrics by correlation magnitude $r$. The primary determinant is cell transition frequency, formalized as the \textbf{cell-switch rate}. For a given subscriber trajectory, the cell-switch rate measures the proportion of consecutive event records registering a change in cell identifier. Stationary subscribers exhibit cell-switch rates near 0\%, whereas highly mobile subscribers approach 100\%. Cell-switch rate demonstrates a strong inverse correlation with prediction accuracy ($r = -0.971$ across 1,173 raw log subscribers and $r = -0.773$ across 100 held-out subscribers in \Dmob{}): \emph{higher cell transition frequency strongly correlates with lower prediction accuracy}.

\textbf{Figure~\ref{fig:concentration}} also demonstrates that alternative hypotheses exhibit weak correlation ($|r| < 0.24$): unique cell count, prediction instance count, and trajectory length show minimal correlation with accuracy. Consequently, trajectory repetitive pattern frequency remains the dominant predictability factor.

\subsection{Quintile Stratification Analysis}

% F7 -- Accuracy is dictated by how often the user changes cell.
\begin{figure}[t]
\centering
\replegend{\swatch{clbase}~persistence baseline \qquad \swatch{clmodel}~model}
\fitwidth{%
\begin{tikzpicture}
\begin{groupplot}[
  group style={group size=2 by 1, horizontal sep=1.5cm, ylabels at=edge left},
  repbar, width=0.47\linewidth, height=5.2cm,
  ymin=0, ymax=125, ytick={0,20,...,100},
  ylabel={top-1 accuracy (\%)},
  /pgf/bar width=7pt, xtick=data,
  x tick label style={font=\scriptsize,rotate=35,anchor=east},
  title style={font=\footnotesize,yshift=1pt},
]
\nextgroupplot[title={(a) \srawbig, five equal groups},
  symbolic x coords={qa,qb,qc,qd,qe},
  xticklabels={0\%,0--21\%,21--38\%,38--60\%,60--100\%},
  xlabel={share of steps where the cell changes}]
\addplot[fill=clbase,draw=none] coordinates {(qa,100.0) (qb,88.4) (qc,71.7) (qd,52.8) (qe,26.6)};
\addplot[fill=clmodel,draw=none] coordinates {(qa,100.0) (qb,87.6) (qc,71.6) (qd,56.5) (qe,34.4)};
\node[font=\scriptsize,text=clgain,anchor=west,rotate=90,yshift=7pt] at (axis cs:qd,74) {$+3.7\pt$};
\node[font=\scriptsize,text=clgain,anchor=west,rotate=90,yshift=7pt] at (axis cs:qe,52) {$+7.8\pt$};
%
\nextgroupplot[title={(b) \smob, four equal groups},
  symbolic x coords={pa,pb,pc,pd},
  xticklabels={0--25\%,25--50\%,50--75\%,75--100\%},
  xlabel={quartile of movement rate}]
\addplot[fill=clbase,draw=none] coordinates {(pa,87.7) (pb,79.6) (pc,66.7) (pd,42.8)};
\addplot[fill=clmodel,draw=none] coordinates {(pa,80.6) (pb,74.3) (pc,63.3) (pd,43.4)};
\node[font=\scriptsize,text=clloss,anchor=west,rotate=90,yshift=7pt] at (axis cs:pa,105) {$-7.1\pt$};
\node[font=\scriptsize,text=clloss,anchor=west,rotate=90,yshift=7pt] at (axis cs:pb,97) {$-5.3\pt$};
\node[font=\scriptsize,text=clloss,anchor=west,rotate=90,yshift=7pt] at (axis cs:pc,84) {$-3.4\pt$};
\node[font=\scriptsize,text=clgain,anchor=west,rotate=90,yshift=7pt] at (axis cs:pd,61) {$+0.6\pt$};
\end{groupplot}
\end{tikzpicture}}
\caption{Accuracy by cell-switch rate}
\label{fig:switch-rate}
\figtext{%
Accuracy is set by how repetitive the user already is, and the model's
contribution runs the other way. In both panels the users are sorted by their
cell-switch rate, the share of consecutive steps at which the cell changes,
and cut into equal groups, five on the left and four on the right, so
volatility increases towards the right; each group carries the persistence
baseline in grey-blue and the model in dark blue, scored on identical
positions, with the gap annotated above. (a)~On the unfiltered day the fifth
of users who never change cell are predicted at 100\,\% by both, because for
them the task is to notice that nothing is happening, and the most volatile
fifth plateaus at 34.4\,\%. That is a 65-point spread produced entirely by who
the users are, on a single trained model, with no hyperparameter involved. The
model overtakes the baseline only in the two hardest groups, by $+3.7$ and
$+7.8\pt$, and those two hold 18.5\,\% of the predictions between them, the
most volatile one alone holding 5.6\,\%: eight points on one twentieth of the
volume moves the aggregate by less than half a point. (b)~The same ordering on
the 100 held-out mobile users of \smob{}, and less comfortably. The model
\emph{loses} to the baseline in the three easier quartiles ($-7.1$, $-5.3$,
$-3.4\pt$) and only draws level in the most volatile one ($+0.6\pt$), which is
precisely how a model can be genuinely useful on hard users and still be worse
than a one-line rule overall. Baselines are recomputed per user, on the same
positions the model is scored on.}
\end{figure}


\textbf{Figure~\ref{fig:switch-rate}} validates this dependency without relying exclusively on scalar correlation coefficients by stratifying subscribers into five equal quintile cohorts based on cell-switch rate.

To verify that individual difficulty reflects intrinsic subscriber behavior rather than evaluation prefix length artifacts, cross-context stability was evaluated. The correlation between individual subscriber accuracy at 30\% context fraction versus 80\% context fraction reaches $r = 0.77$ (\textbf{Table~\ref{tab:context-collapse}} in Chapter~\ref{ch:time}), confirming that predictability is an intrinsic subscriber trajectory property.

\subsection{Performance Breakdown across Restlessness Segments}

\textbf{\panelref{fig:switch-rate}{a}} highlights model performance across cell-switch quintiles. Across the three lowest cell-switch rate quintiles of \textbf{\panelref{fig:switch-rate}{a}} (low mobility), fine-tuned model accuracy and persistence baseline performance are virtually identical. Model accuracy gains manifest exclusively in the two highest cell-switch rate quintiles of \textbf{\panelref{fig:switch-rate}{a}} (high mobility), exceeding baseline performance by +3.7 and +7.8 percentage points, respectively.

However, high cell-switch subscribers account for only 18.5\% of total evaluated prediction instances in \textbf{\panelref{fig:switch-rate}{a}}, with the highest cell-switch quintile representing 5.6\% of total volume. Achieving a +7.8 percentage point improvement on 5.6\% of total evaluation volume yields a negligible shift in aggregate population accuracy. This establishes a key reporting convention: \emph{aggregate population accuracy averages distinct mobility cohorts; reporting segment-specific accuracy relative to baseline benchmarks is required to capture true model contributions.}

\section{Impact of User Activity Volume on Aggregate Metrics}

Evaluating prediction volume distribution reveals significant inequality across subscribers. Partitioning 1,173 subscribers into quintiles based on event volume demonstrates that the highest-volume decile (117 subscribers) contributes \textbf{58.6\%} of the 54,173 total evaluated prediction instances. Consequently, unweighted aggregate accuracy metrics are heavily dominated by this sub-population.

However, event volume does not correlate strongly with prediction accuracy ($r = +0.154$). Higher event volume corresponds to slightly lower cell-switch rates ($r = -0.150$), explaining the minor upward accuracy trend in high-volume cohorts.

\section{Analysis of Cell-Level vs. Mast-Level Prediction Errors}

Analyzing error cases at individual subscriber resolution reveals that a portion of prediction errors stems from cell-level vs. base-station-level definitions.

% T4 -- The seven worst users, and what actually separates them (slide 28).
\begin{table}[!htbp]
\centering\small
\caption{The seven worst-scoring people, and the reason each scores badly}
\label{tab:worst-users}
\begin{tabular}{@{}lrrrrl@{}}
\toprule
\textbf{\texttt{user\_id}} & \textbf{Moments} & \textbf{Ceiling} & \textbf{Rule} & \textbf{Model} & \textbf{Which situation} \\
\midrule
1049149 & 3  & 0\,\%   & 0\,\%  & 0\,\%    & \multirow{4}{*}{\begin{tabular}{@{}l@{}}\textbf{A}, two to twelve moments each:\\ too little material to judge a model\end{tabular}} \\
1291964 & 3  & 0\,\%   & 0\,\%  & 0\,\%    & \\
1344347 & 8  & 25\,\%  & 12\,\% & 12.5\,\% & \\
929903  & 12 & 75\,\%  & 8\,\%  & 8.3\,\%  & \\
\addlinespace
1161518 & 9  & 89\,\%  & 67\,\% & 11.1\,\% & \textbf{B}, a real change of behaviour \\
\addlinespace
828479  & 2  & 100\,\% & 50\,\% & 0\,\%    & \multirow{2}{*}{\begin{tabular}{@{}l@{}}\textbf{C}, two names, one mast\end{tabular}} \\
1067194 & 3  & 100\,\% & 67\,\% & 0\,\%    & \\
\bottomrule
\end{tabular}

\figsource{\nbllm}{train\_cellid\_llm.ipynb}{the per-user error listing}
\end{table}


\textbf{Table~\ref{tab:worst-users}} details the seven lowest-scoring subscriber trajectories from a representative evaluation run. Subgroups at the bottom of \textbf{Table~\ref{tab:worst-users}} oscillate between adjacent directional sector antennas mounted on the same physical base station mast. On these trajectories, the persistence baseline achieves 50\%--67\% accuracy by repeating sector IDs, whereas the model attempts to forecast specific sector switching, incurring prediction errors despite zero net spatial displacement.

Maintaining granular \cid{cell\_id} target prediction fulfills the cell-level forecasting mandate, even though a fraction of prediction errors represents intra-site sector handovers rather than geographical mobility.

\section{Summary of Key Findings and Temporal Hypothesis}

Three empirical findings summarize this chapter:
1. Individual trajectory predictability is dictated primarily by cell-switch frequency (repetitiveness).
2. Fine-tuned language model accuracy gains concentrate within high mobility cohorts representing a minor fraction of total evaluation events.
3. Intra-site sector transitions account for a systematic fraction of cell-level prediction errors.

Because scaling general capacity does not alter low-mobility baseline bounds, improving high-mobility transition forecasting requires temporal conditioning signals (hour of day) to anticipate transitions.

\section{Conclusion}
\label{sec:retro-user}

This chapter demonstrated that trajectory predictability is primarily governed by subscriber mobility characteristics rather than model parameter scale. Cell-switch rate represents the single strongest predictor of accuracy ($r = -0.773$), whereas trajectory length and unique cell counts exhibit minimal correlation ($|r| < 0.24$).

Model accuracy gains over persistence baselines concentrate within high mobility cohorts (+7.8 percentage points in the highest cell-switch quintile of \textbf{\panelref{fig:switch-rate}{a}}). Because high mobility subscribers generate a minority of total event volume, aggregate population accuracy shifts marginally. These findings establish the necessity of reporting cohort-stratified accuracy deltas throughout mobility prediction research.

